\section{The near-quadratic sampling problem}\label{sec:preliminaries}

The conditional density in the introduction is a Gaussian multiplied
by a factor with a small Hessian. We now formulate the sampling
problem for that larger class of perturbations. The center will be
arbitrary: later calls must work at centers supplied by earlier
random computations, without a bound on the gradient there. We state
the required uniform guarantee first, then prove the comparison and
smoothing facts used throughout its construction.

\subsection{Notation and the auxiliary theorem}

Vectors have Euclidean norm, and matrices have its induced operator
norm unless otherwise indicated. For a random vector or tensor,
\(\norm{X}_2=(\E|X|^2)^{1/2}\), with the chosen pointwise norm;
the notation \(\norm{X}_{\HS,2}\) specifies the Frobenius norm.
The distance \(\Wtwo\) is quadratic Wasserstein distance on
Euclidean space. Total variation is
\(\norm{\mu-\eta}_{\TV}=\sup_B|\mu(B)-\eta(B)|\).
All probability laws and oracle procedures are measurable.

Fix \(0<e<1\), and put \(\delta=d^{-e}\). Let
\(F:\R^d\to\R\) be \(C^2\), with gradient \(g=\nabla F\)
of Lipschitz constant at most the declared bound
\(\lambda=\delta\). Values of \(F\) need not be available.
For \(x\in\R^d\) and \(0<r\le1\), define
\begin{equation}\label{eq:primitive-law}
 \nu_{x,r}(dz)=
 \frac{\exp(-|z|^2/2-F(x+rz))\,dz}
 {\int_{\R^d}\exp(-|v|^2/2-F(x+rv))\,dv},
 \qquad f(z)=g(x+rz),\qquad u=\E_{\nu_{x,r}}f.
\end{equation}
We often write \(\nu=\nu_{x,r}\), and use
\begin{equation}\label{eq:primitive-scales}
 L=\lambda r,\qquad l=\lambda r^2=rL,\qquad
 D(x)=\sqrt d+|g(x)|.
\end{equation}
Here \(x\) is the actual input center; it need not be a minimizer
of \(F\). The positive number \(\lambda\) is a declared upper bound,
so the notation also covers affine or constant \(F\).
The full potential in \eqref{eq:primitive-law} has Hessian between
\((1-l)I\) and \((1+l)I\). Since \(l\le\delta<1\) for \(d\ge2\),
the normalizing integral is finite, even when \(F\) is not convex.
For the quantitative construction we take \(d\) large enough that
\(l\le1/2\).

Throughout the construction, \(\polylog(d)\) denotes a factor
bounded by \(C(1+\log d)^C\). Its constants may depend on the
fixed parameters \(e\), accuracy orders, and other fixed orders
introduced below, but never on \(F,x,r\), or the oracle responses.
Different occurrences may have different constants. Every
sufficiently-large-dimension threshold has the same uniformity.

\begin{theorem}[Noisy near-quadratic sampling]\label{thm:primitive}
Fix \(0<e<1\), an integer \(K_0\ge2\), and \(t>0\) satisfying
\begin{equation}\label{eq:t-choice}
                 40(K_0+2)t<\frac14.
\end{equation}
There are constants and a dimension threshold depending only on
\(e,K_0,t\) with the following property. For every integer \(d\)
above that threshold, put \(\delta=d^{-e}\). One measurable randomized
oracle procedure works for every \(F\in C^2(\R^d)\) with
\(\Lip(\nabla F)\le\delta\), every \(x\in\R^d\), and every
\(0<r\le1\). Given evaluation access to \(g=\nabla F\), it
uses at most \(\polylog(d)\delta^{-t}\) evaluations of
\(g\) on every execution and returns
\[
                  X=X_{\mathrm{pre}}+\tfrac12G,
\]
where \(G\sim N(0,I_d)\) is independent of \(X_{\mathrm{pre}}\) and of all
queries, and
\begin{equation}\label{eq:primitive-pre-noise}
 \Wtwo\!\left(\mathcal L(X_{\mathrm{pre}}),
       \nu_{x,r}*N(0,\tfrac34I_d)\right)
       \le \polylog(d)\delta^{K_0}D(x).
\end{equation}
Consequently both
\[
 \Wtwo\!\left(\mathcal L(X),\nu_{x,r}*N(0,I_d)\right)
 \quad\hbox{and}\quad
 \norm{\mathcal L(X)-\nu_{x,r}*N(0,I_d)}_{\TV}
\]
are at most \(\polylog(d)\delta^{K_0}D(x)\).
All constants in these bounds are uniform in \(F,x,r\).
The procedure's scalar parameters, shapes of its computations,
and query cap are independent of \(x\) and of the primitive.
It uses no evaluations of \(F\) itself.
\end{theorem}

The proof occupies Sections~\ref{sec:jets}--\ref{sec:compiler}.
The independent Gaussian is reserved until all queries have been made.
Adding the same reserve to a coupling preserves its transportation
error; smoothing by that reserve gives the total-variation conclusion
through Lemma~\ref{lem:coupling-smoothing} below. The target here
contains a full independent Gaussian convolution.
Section~\ref{sec:sampling} will remove that auxiliary noise before
returning a sample for the original problem.

\subsection{Convexity, moments, and comparisons}

We recall the elementary diffusion facts needed for the proof,
including their regularity requirements. They are familiar forms
of the curvature method for functional inequalities
\citep{bakry2014analysis}; we give the argument used here.

\begin{lemma}[Uniformly convex laws]\label{lem:convex-preliminaries}
Let \(H\in C^2(\R^n)\) satisfy
\(mI\preceq\nabla^2H\preceq MI\), with \(0<m\le M<\infty\).
Write \(\mu(dx)\propto e^{-H(x)}dx\), and let \(z_0\) be
its unique mode. Then
\[
 \E_\mu|X-z_0|^2\le n/m,\qquad
 \Var_\mu h\le m^{-1}\E_\mu|\nabla h|^2.
\]
The latter applies, in particular, to \(C^1\) functions whose
value and gradient have polynomial growth.
The diffusion
\[
                  dX_t=-\nabla H(X_t)\,dt+\sqrt2\,dB_t
\]
has invariant law \(\mu\), and synchronous solutions contract
in distance by \(e^{-mt}\).

If \(\bar H\) is another strongly convex \(C^2\) potential
with bounded Hessian and law \(\eta\), then
\begin{equation}\label{eq:synchronous-comparison}
 \Wtwo(\mu,\eta)
 \le \frac1m
       \left(\E_\eta|\nabla H-\nabla\bar H|^2\right)^{1/2}.
\end{equation}
\end{lemma}

\begin{proof}
Strong convexity implies Gaussian tails and existence of a
unique minimizer. Integration by parts gives
\(\E_\mu (X-z_0)\cdot\nabla H(X)=n\); strong monotonicity
of \(\nabla H\) gives the moment bound.
The globally Lipschitz drift gives a nonexplosive diffusion.
For synchronous solutions, differentiating their squared
distance and using strong monotonicity gives the contraction.

First suppose \(H\) is smooth and has bounded derivatives of
orders at least two. For a smooth compactly supported \(h\),
let \(P_th(x)=\E[h(X_t)\mid X_0=x]\).
Differentiating the synchronous flow in the initial state
gives bounded first and second derivatives on bounded time
intervals. It\^o's formula and the semigroup identity yield
\[
 \partial_tP_th=P_t\mathcal Ah=\mathcal AP_th,
 \qquad \mathcal A=\Delta-\nabla H\cdot\nabla.
\]
The equality with \(\mathcal AP_th\) also follows by
differentiating \(P_{t+s}h=P_s(P_th)\) at \(s=0\).
Its at-most-linear growth is integrable against \(\mu\).
Integration by parts gives
\(\int\mathcal AP_th\,d\mu=0\), hence invariance.
Coupling to an invariant starting point and using contraction
shows \(P_th\to\mu h\) in \(L^2(\mu)\). Therefore
\[
 \Var_\mu h=2\int_0^\infty
                \int|\nabla P_th|^2\,d\mu\,dt.
\]
The flow derivative has operator norm at most \(e^{-mt}\).
Jensen's inequality and invariance thus bound the inner
integral by \(e^{-2mt}\int|\nabla h|^2d\mu\), proving
the Poincar\'e inequality.

For \(C^2\) potentials, convolve with smooth compactly
supported mollifiers of radii tending to zero. Both Hessian
bounds are preserved; all higher derivatives are bounded
for each mollification. The normalized laws converge in
total variation by the uniform quadratic tail bounds
(subtract the potential's value at zero).
The gradients converge uniformly since the original
gradient is globally Lipschitz. Synchronous solutions of
the mollified and original diffusions consequently converge
uniformly in their common starting point on each bounded
time interval, by Gronwall's inequality. Testing invariance
against bounded Lipschitz functions passes it to the limit.
Poincar\'e first passes on smooth compact tests, then by
cutoff and approximation to the stated class.

For \eqref{eq:synchronous-comparison}, start the two
diffusions with their respective invariant marginals,
with any initial coupling, and use common Brownian motion.
Writing \(a(t)=(\E|X_t-Y_t|^2)^{1/2}\) and
\(\beta=(\E_\eta|\nabla H-\nabla\bar H|^2)^{1/2}\),
strong monotonicity and Cauchy--Schwarz give, in the usual
upper-derivative sense,
\[
                         a'(t)\le-ma(t)+\beta.
\]
The marginals remain fixed, so \(\Wtwo(\mu,\eta)\le a(t)\).
Letting \(t\) tend to infinity proves the claim.
\end{proof}

These estimates apply to the auxiliary laws without imposing convexity
on \(F\). For \(\nu_{x,r}\), the mode equation is
\(z_0=-r g(x+rz_0)\). In particular, for \(l\le1/2\),
\begin{equation}\label{eq:primitive-moments}
 |z_0|\le\frac{r|g(x)|}{1-l},\qquad
 \bigl(\E_\nu|Z|^2\bigr)^{1/2}
      \le C(\sqrt d+r|g(x)|),\qquad |u|\le CD(x).
\end{equation}
Indeed, \(|g(x+rz_0)-g(x)|\le L|z_0|\) proves the first bound,
and the lemma gives \(\norm{Z-z_0}_2\le\sqrt{d/(1-l)}\).
The triangle inequality and
\(|u|\le|g(x)|+L\norm{Z}_2\) give the remaining bounds.
Thus \(D(x)\) controls the required moments uniformly over all
centers, including those with very large gradients.

\subsection{Coupling and smoothing conventions}

Errors of a deterministic map evaluated on a specified
input law will be measured in root mean square. Errors of
randomized outputs will be measured in \(\Wtwo\). The
following rules allow us to keep these two notions separate.

\begin{lemma}[Conditional comparisons and smoothing]\label{lem:coupling-smoothing}
Let \((\mu_z)\) and \((\eta_z)\) be probability kernels on
finite-dimensional Euclidean spaces, and let \(Z\) have
law \(\zeta\). Whenever the relevant second moments are finite,
\[
 \Wtwo^2\!\left(\int\mu_z\,d\zeta(z),\int\eta_z\,d\zeta(z)\right)
       \le\int\Wtwo^2(\mu_z,\eta_z)\,d\zeta(z).
\]
For finite families, independent sums satisfy
\[
 \Wtwo\!\left(\mathcal L\Bigl(\sum_i a_iX_i\Bigr),
             \mathcal L\Bigl(\sum_i a_iY_i\Bigr)\right)
 \le\sum_i|a_i|\Wtwo(\mathcal L(X_i),\mathcal L(Y_i)),
\]
where independence is required within each of the two
families, not between them.
For \(s>0\),
\begin{equation}\label{eq:gaussian-smoothing}
 \norm{\mu*N(0,s^2I)-\eta*N(0,s^2I)}_{\TV}
                       \le s^{-1}\Wtwo(\mu,\eta).
\end{equation}
Moreover, applying a common probability kernel cannot
increase total variation, and
\[
 \left\|\int\mu_z\,d\zeta(z)-\int\eta_z\,d\zeta(z)\right\|_{\TV}
       \le\int\norm{\mu_z-\eta_z}_{\TV}\,d\zeta(z).
\]
\end{lemma}

\begin{proof}
Use conditional near-optimal couplings for the first
inequality and their independent product for the second,
followed by the triangle inequality in \(L^2\).
Measurability of such conditional comparisons can also be
seen without an exact optimal-kernel selection: quantize
each output law on a finite grid inside a box, sending
the complement to zero. A minimizing plan on the finite
sets can be chosen measurably by a fixed ordering of the
finitely many basic feasible solutions. Lift it using
the conditional measures in the cells, and refine the
grids and boxes. Finite unconditional second moments
and the triangle inequality give the asserted limit.

For deterministic centers \(x,y\), the relative entropy
of \(N(x,s^2I)\) from \(N(y,s^2I)\) is
\(|x-y|^2/(2s^2)\). Pinsker's inequality gives TV at most
\(|x-y|/(2s)\). Integrate this estimate over any coupling
and apply Cauchy--Schwarz to obtain
\eqref{eq:gaussian-smoothing}, with room in its constant.
The total-variation rules follow directly by testing
measurable functions taking values in \([0,1]\).
\end{proof}

We will also repeatedly combine deterministic root-mean-square
errors with absolute coefficient sums. This uses the
ordinary \(L^2\) triangle inequality and does not require
independence.
